\section{Discussion, Limitations, and Conclusion}
\label{sec:discussion}

\subsection{What the evidence establishes}

Three findings form the paper's causal spine. First, the bandwidth of the training target is a genuine causal control of spectral fabrication, not a correlate: changing only that choice, from an identical initialization and with identical additional optimization, cuts measured out-of-band energy to about one third of the matched control, 6/6, with the no-reference forward-consistency endpoint ruling out a smoothness-degradation explanation. Second, fabrication is not free of capacity: in the two tested widths it is the larger network that fabricates more, and that admits stronger suppression---capacity that bandwidth-matched targets redirect rather than remove. Third, the linear Wiener baseline shows what the trade-off landscape looks like from the outside: fabrication can be pushed below the learned range, but only by accepting a roughly $1.4\times$ (about $40\%$) worse forward fit. The practical statement for practitioners is deliberately modest and operational: audit your training target's bandwidth against your system's measured support before auditing your network.

\subsection{Limitations}

The causal claims are established under simulated protocols: a shift-invariant speckle model (plus the position-dependent robustness variant), a single noise model, $64\times64$ images, two network widths, three seeds, and digit objects. Out-of-band energy is a fabrication proxy anchored to a measured physical support, not a perceptual or semantic metric. The capacity result covers two widths and is stated as an observed trend. The distance-law negative control relies on three metric choices and twelve diffuser families. Real-data evidence is preliminary: system parameters and pairing statistics are quantified, but no quantitative real-data reconstruction claim is made, and the failed registration approaches delimit rather than solve the forward-calibration problem. The fine-tune comparison, though compute-matched, involves a single fine-tune length per setting in the main arm (the length sweep in Sec.~\ref{sec:svpsf} addresses stability, not exhaustiveness).

\subsection{Future work}

Two directions follow directly. (i) \emph{Objects beyond digits and one real forward-calibration}: extending objects to natural-image statistics and building the forward-calibration stage that the real-data section identifies would let the protocol cross from simulation to measured systems. (ii) \emph{Bandwidth as a deployment dial}: the intervention costs nothing and is compatible with any training pipeline; characterizing it under temporal integration regimes and multi-wavelength acquisition would test whether measured support should be treated as a first-class hyperparameter of learned reconstruction.

\subsection{Conclusion}

We isolated the bandwidth of the training target as a causal, compute-matched control of spectral fabrication in learned speckle reconstruction. Bandwidth-matched fine-tuning reduces out-of-band fabrication to $0.34 \pm 0.11$ (matched) and $0.20 \pm 0.08$ (half) of an identically optimized control (6/6 units, sign test $p \approx 0.016$) while leaving the no-reference forward-consistency endpoint nearly unchanged ($+0.026$ on a $0.36$--$0.40$ baseline); raw fabrication and suppression strength both grow with capacity in the tested range; and a Wiener anchor shows the control operates within a trade-off landscape where the spectrally conservative linear baseline remains under-fitted by about $40\%$ in relative forward error. A cross-diffuser analysis finds no class boundary, and a real-data assessment quantifies system parameters while marking the forward-calibration requirement as the boundary of the simulation-based claim. Fabrication in learned speckle imaging is, in part, a choice made in the training data---and it can therefore be audited, controlled, and reported as one.
